\documentclass{article}
\usepackage[utf8]{inputenc}
\usepackage{geometry}
\geometry{a4paper, margin=1in}

\title{Respuestas - Evaluación Semana 1 (TXT/SOA)}
\author{José}
\date{\today}

\begin{document}

\maketitle

\section*{Soluciones al Cuestionario}

\begin{enumerate}
    \item \textbf{Información de un registro SOA:}
    Un registro \textbf{SOA} (Start of Authority) proporciona información administrativa esencial sobre una zona DNS. Incluye datos como:
    \begin{itemize}
        \item El servidor de nombres primario para la zona.
        \item El correo electrónico del administrador de la zona.
        \item El número de serie de la zona (que indica si ha habido cambios).
        \item Temporizadores relacionados con la actualización y expiración de la zona.
    \end{itemize}

    \item \textbf{Registros TXT y SPF:}
    Un registro \textbf{SPF} (Sender Policy Framework) es un tipo de registro \textbf{TXT} que define una lista de servidores de correo autorizados para enviar correos electrónicos en nombre de un dominio. Ayuda a mitigar el problema de \textbf{suplantación de correo electrónico (email spoofing)} y el spam, ya que permite a los servidores receptores verificar si un correo que dice provenir de un dominio (ej. \texttt{@miempresa.com}) fue realmente enviado desde un servidor autorizado por \texttt{miempresa.com}.

    \item \textbf{Excepción en \texttt{dnspython}:}
    Cuando se realiza una consulta para un tipo de registro válido (como 'TXT' o 'SOA') pero el dominio consultado no tiene ningún registro de ese tipo configurado, la biblioteca \texttt{dnspython} lanza la excepción \texttt{dns.resolver.NoAnswer}. Esta excepción indica que la consulta fue exitosa, pero la respuesta del servidor DNS estaba vacía.

\end{enumerate}

\end{document}